Let's put the change of variables to work.

\begin{eg}[Computing the area of the disc]
	We consider the disk $B_{1} = B_{1}(0) \subset \mathbb{R}^{2}$ and wish to compute $\mu_{2}(B_{1}) = \pi$. We define
	\begin{align*}
		\Psi(y_{1}, y_{2}) = \begin{pmatrix} y_{1} \cos y_{2} \\ y_{1} \sin y_{2} \end{pmatrix}
	\end{align*}
	defined on $D = (0,1) \times (0, 2\pi)$ and mapping into $B_{1}$. And notice that although $B_{1} \neq \Psi(D)$, the part of $B_{1}$ that is not covered by $\Psi$ has measure zero, hence $\mu_{2}(B_{1}) = \mu_{2}(\Psi(D))$. But by Lemma \ref{lem:int-const-func},
	\begin{align*}
		\mu_{2}(\Psi(D)) = \int_{\Psi(D)} 1 \, dx = \int_{D} \left| \det J \Psi(y) \right| \, dy.
	\end{align*}
	And actually, one computes $\det J \Psi(y_{1}, y_{2}) = y_{1} > 0$, so we're looking at
	\begin{align*}
		\int_{(0,1) \times (0,2\pi)} y_{1} \, dy.
	\end{align*}
	But, using Lemma \ref{lem:box-volume}, the volume of this (consider a drawing of the hypograph) is $\frac{1}{2} \cdot \mu_{3}((0,1) \times (0,2\pi) \times (0,1)) = \frac{1}{2} \cdot 2\pi = \pi$. And so we know the area of the disk.
\end{eg}

\begin{eg}[Computing the volume of the sphere]
	We wish to compute $\mu_{3}(B_{1}) = \frac{4}{3} \pi$ for $B_{1} \subset \mathbb{R}^{3}$. For that, we have the polar coordinates again
	\begin{align*}
		\Psi(y_{1}, y_{2}, y_{3}) =
		\begin{pmatrix}
			y_{1} \cos y_{2} \cos y_{3} \\
			y_{1} \sin y_{2} \cos y_{3} \\
			y_{1} \sin y_{3}
		\end{pmatrix}
	\end{align*}
	with domain $D = (0,1) \times (0,2\pi) \times (-\tfrac{\pi}{2}, \tfrac{\pi}{2})$. One computes $\left| \det J \Psi \right| = y_{1}^{2} \cos y_{3}$ an so we have to compute
	\begin{align*}
		\mu_{3}(B_{1}) = \mu_{3}(\Psi(D)) = \int_{(0,1) \times (0, 2\pi) \times (-\tfrac{\pi}{2}, \tfrac{\pi}{2})} y_{1}^{2} \cos y_{3} \, dy.
	\end{align*}
	Now, to compute this integral, we have to introduce the slicing formula.
\end{eg}

\section{Cavalieri's principle and Fubini's theorem}

% drawing of cavalieri

The Cavalieri principle in 3D says roughly, that for any measurable set $A \subset \mathbb{R}^{n}$, you can compute the volume of $A$ by partitioning it into very thin slices, and to obtain the total volume, you just have to sum over all the little slices.

For the following proofs, split the notation of points in $\mathbb{R}^{n}$ as points in $\mathbb{R}^{k} \times \mathbb{R}^{n-k}$. So writing $(x,y) \in \mathbb{R}^{n}$ implies that $x \in \mathbb{R}^{k}$ and $y \in \mathbb{R}^{n-k}$.

\begin{theorem}[Cavalieri]\label{thm:cavalieri}
	Let $A \subset \mathbb{R}^{n}$ be Jordan measurable, $A \subset (-C, C)^{n}$. Define the slices for $x \in (-C, C)^{k}$ as
	\begin{align*}
		A_{x} = \left\{ y \in \mathbb{R}^{n-k} \mid (x,y) \in A \right\} \subset (-C, C)^{n-k}
	\end{align*}
	and their inner and outer volumes in $\mathbb{R}^{n-k}$
	\begin{align*}
		\overline{ \varphi } (x) = \mu_{n-k, \text{out}}(A_{x}), \quad \underline{ \varphi}(x) = \mu_{n-k, \text{in}}(A_{x}).
	\end{align*}
	Then $\overline{ \varphi }, \underline{\varphi}: (-C, C)^{k} \to \mathbb{R}$ are Riemann integrable and
	\begin{align*}
		\int_{(-C, C)^{k}} \overline{ \varphi }(x)  \, dx = \int_{(-C, C)^{k}} \underline{ \varphi }(x) \, dx = \mu_{n}(A).
	\end{align*}
\end{theorem}
\begin{proof}
	Given a dyadic cube of pixel size $2^{-p}$ in $\mathbb{R}^{n}$, we write
	\begin{align*}
		Q_{p}^{n}(z) = z + 2^{-p}[0,1)^{n}.
	\end{align*}
	With $z = (x,y) \in \mathbb{R}^{k} \times \mathbb{R}^{n-k}$, we can rewrite the cube as
	\begin{align*}
		Q_{p}^{n}(z) = Q_{p}^{k}(x) \times Q_{p}^{n-k}(y).
	\end{align*}
	Since $A$ is measurable, given $\epsilon > 0$, let $F, G$ dyadic sets of size $2^{-p}$ such that $F \subset A \subset G$ and $\mu(G) - \mu(F) < \epsilon$. For $x \in (-C, C)^{k}$ define
	\begin{align*}
		G_{x} = \left\{ y \in \mathbb{R}^{n-k} \mid (x,y) \in G \right\}, \quad F_{x} = \left\{ y \in \mathbb{R}^{n-k} \mid (x,y) \in F \right\}.
	\end{align*}
	Now observe that $F_{x} \subset A_{x} \subset G_{x}$ and $F_{x} = F_{x'}$ for all $x, x' \in Q^{k}_{p}(a)$, because the cubes of $F$ can be "factored" as $Q_{p}^{n}((x,y)) = Q_{p}^{k}(x) \times Q_{p}^{n-k}(y)$. Same holds for $G_{x}$.

	So $f(x) = \mu(F_{x})$ and $g(x) = \mu(G_{x})$ are both constant on each cube $Q_{p}^{k}(a)$, hence dyadic step functions for which it holds that $f(x) \leq \underline{ \varphi }(x) \leq \overline{ \varphi }(x) \leq g(x)$. Now, each cube $Q_{p}^{k}(a) \times Q_{p}^{n-k}(b)$ contributes function value $2^{-p(n-k)}$ times domain base length $2^{-pk}$, so 
	\begin{align*}
		f(x) = \mu(F_{x}) = 2^{-pn} \cdot \text{\# of dyadic cubes in $F_x$}.
	\end{align*}
    and if we take the integral over all possible $Q_{p}^{k}(a)$ cubes, 
	\begin{align*}
		\int_{\mathbb{R}^{k}} f(x) \, dx = 2^{-pn} \cdot \text{\# of dyadic cubes in $F$} = \mu(F),
	\end{align*}
	and analogously,
	\begin{align*}
		\int_{\mathbb{R}^{k}} g(x) \, dx = 2^{-pn} \cdot \text{\# of dyadic cubes in $G$} = \mu(G).
	\end{align*}
	Hence
	\begin{align*}
		\mu(F) = \int f(x) \, dx \leq \underline{ \int } \underline{ \varphi }(x) \, dx \leq \overline{ \int } \overline{ \varphi } (x) \, dx \leq \int g(x) \, dx = \mu(G).
	\end{align*}
	But since $\mu(G) - \mu(F) < \epsilon$, $\mu(F) \leq \mu(A) \leq \mu(G)$ and $\epsilon > 0$ was arbitrary, we get the entire equality.
\end{proof}



